\section{The extremal contradiction}\label{sec:incompatibility}

We use Lemma~\ref{lem:pruning} to rule out the graph supplied by
Proposition~\ref{prop:counterexample-reduction}.  The proof converts the strict
subset inequalities into a strict increase of an extremal quantity associated
with two families of ordered pairs.

\begin{proposition}\label{prop:no-concavity}
There is no nonempty finite oriented graph $G$ in which every vertex has an
in-neighbor and whose one-step image operator
\[
 F(U)=\{y\in V(G):x\to y\text{ for some }x\in U\}
\]
satisfies
\begin{equation}\label{eq:impossible-concavity}
 \abs{U}+\abs{F^2(U)}<2\abs{F(U)}
 \qquad(\varnothing\ne U\subsetneq V(G)),
\end{equation}
where $F^2=F\circ F$.
\end{proposition}

\begin{proof}
Suppose that such a graph exists, and put $X=V(G)$.  For $S\subseteq X\times X$,
define the two coordinate image operators by
\[
 \begin{aligned}
 F_1S&=\{(i,j): (p,j)\in S\text{ and }p\to i
                     \text{ for some }p\in X\},\\
 F_2S&=\{(i,j): (i,s)\in S\text{ and }s\to j
                     \text{ for some }s\in X\}.
 \end{aligned}
\]
Write $F_1^2=F_1\circ F_1$, $F_2^2=F_2\circ F_2$, and
$\Delta=\{(v,v):v\in X\}$.  Call $P,Q\subseteq X\times X$ \emph{compatible}
if $F_1P\cap F_2Q=\varnothing$.  Equivalently, no left member $(p,j)\in P$
and right member $(i,s)\in Q$ satisfy both $p\to i$ and $s\to j$; this is
precisely the conflict relation in Lemma~\ref{lem:pruning}.

The pair $P=Q=\Delta$ is compatible: a point $(i,j)$ in both images would
require $j\to i$ and $i\to j$, contrary to orientation.  Thus we may choose,
among all compatible pairs with $\Delta\subseteq P\cap Q$, one maximizing
$M+d$, where
\[
 M=\abs{P}+\abs{Q},\qquad
 d=\abs{X}^2-\abs{F_1P}-\abs{F_2Q}.
\]
Such a maximum exists because $X$ is finite.  Compatibility makes $d$ the
number of points outside $F_1P\cup F_2Q$.  Including these uncovered points
in the objective lets the points left uncovered by pruning offset the
corresponding part of the deletion cost.

Every column $P_j=\{p:(p,j)\in P\}$ is nonempty, since $j\in P_j$.
It is also proper.  Indeed, positive indegrees give $F(X)=X$, so $P_j=X$
would imply $X\times\{j\}\subseteq F_1P$.  Choose an in-neighbor $w$ of
$j$.  Then $(w,j)\in F_2\Delta\subseteq F_2Q$, contradicting compatibility.
Likewise every row $Q_i=\{s:(i,s)\in Q\}$ contains $i$.  If $Q_i=X$,
then $\{i\}\times X\subseteq F_2Q$; choosing $w\to i$ gives
$(i,w)\in F_1\Delta\subseteq F_1P$, again a contradiction.
We can therefore sum~\eqref{eq:impossible-concavity} over the columns of $P$
and over the rows of $Q$ to obtain
\begin{equation}\label{eq:strict-fiber-sums}
 \abs{P}+\abs{F_1^2P}<2\abs{F_1P},\qquad
 \abs{Q}+\abs{F_2^2Q}<2\abs{F_2Q}.
\end{equation}

Define new left and right families by taking complements in $X\times X$:
\[
 \mathcal R=(X\times X)\setminus F_2^2Q,
 \qquad
 \mathcal C=(X\times X)\setminus F_1^2P.
\]
The two strict inequalities in~\eqref{eq:strict-fiber-sums} give
\begin{equation}\label{eq:complement-gain}
 \begin{aligned}
 \abs{\mathcal R}+\abs{\mathcal C}
 &=2\abs{X}^2-\abs{F_2^2Q}-\abs{F_1^2P}\\
 &>2\abs{X}^2+M-2\bigl(\abs{F_1P}+\abs{F_2Q}\bigr)
 =M+2d.
 \end{aligned}
\end{equation}
These families need not be compatible.  We will apply the pruning lemma,
first verifying that the diagonal survives and then bounding the cost.

Both families contain $\Delta$.  If $(v,v)\in F_2^2Q$, there is an $s\to v$
with $(v,s)\in F_2Q$.  But $(s,s)\in P$ gives $(v,s)\in F_1P$, a
contradiction.  If $(v,v)\in F_1^2P$, there is a $p\to v$ with
$(p,v)\in F_1P$.  Now $(p,p)\in Q$ gives $(p,v)\in F_2Q$, the same
contradiction.

Moreover, each diagonal member is conflict-free in the new families.
A conflict between the left member $(v,v)$ and a right member
$(i,s)\in\mathcal C$ would give the path $s\to v\to i$.
Starting at $(s,s)\in P$ and expanding its first coordinate twice would
then put $(i,s)$ in $F_1^2P$, contrary to its membership in $\mathcal C$.
A conflict between a left member $(p,j)\in\mathcal R$ and the right member
$(v,v)$ would give $p\to v\to j$.  Expanding the second coordinate of
$(p,p)\in Q$ twice would put $(p,j)$ in $F_2^2Q$, again a contradiction.
The preservation clause of Lemma~\ref{lem:pruning} thus applies to every
diagonal member on both sides.

Let $H,Z\subseteq X\times X$ be the sets associated with
$\mathcal R,\mathcal C$ in that lemma.  Thus each conflict
\[
 (p,j)\in\mathcal R,\quad (i,s)\in\mathcal C,\quad
 p\to i,\quad s\to j
\]
contributes $(p,s)$ to $H$ and $(i,j)$ to $Z$.
We claim that
\begin{equation}\label{eq:old-uncovered-bound}
 H\subseteq (X\times X)\setminus(F_1P\cup F_2Q),
 \qquad\text{and hence}\qquad \abs{H}\le d.
\end{equation}
To see this, fix such a conflict.  If $(p,s)\in F_1P$, choose
$(u,s)\in P$ with $u\to p$.  The path $u\to p\to i$ then puts
$(i,s)$ in $F_1^2P$, contrary to $(i,s)\in\mathcal C$.
If $(p,s)\in F_2Q$, choose $(p,t)\in Q$ with $t\to s$.
The path $t\to s\to j$ puts $(p,j)$ in $F_2^2Q$, contrary to
$(p,j)\in\mathcal R$.  This proves~\eqref{eq:old-uncovered-bound}.

Apply Lemma~\ref{lem:pruning} to obtain conflict-free families
$P'\subseteq\mathcal R$ and $Q'\subseteq\mathcal C$.
They still contain $\Delta$, and their lack of conflicts means
$F_1P'\cap F_2Q'=\varnothing$.  Set
\[
 d'=\abs{X}^2-\abs{F_1P'}-\abs{F_2Q'}.
\]
A point is covered by a kept left member exactly when it belongs to
$F_1P'$, and by a kept right member exactly when it belongs to $F_2Q'$.
Consequently at most $d'$ points of $Z$ are uncovered after pruning.
By Lemma~\ref{lem:pruning} and~\eqref{eq:old-uncovered-bound}, at most $d+d'$ members
were deleted, counting the two sides separately.  Using
\eqref{eq:complement-gain}, we conclude that
\[
 \begin{aligned}
 \abs{P'}+\abs{Q'}+d'
 &\ge \abs{\mathcal R}+\abs{\mathcal C}-(d+d')+d'\\
 &>M+2d-(d+d')+d'=M+d.
 \end{aligned}
\]
This contradicts the choice of $P,Q$ and proves the proposition.
\end{proof}

\begin{proof}[Proof of Theorem~\ref{thm:main}]
If the theorem were false, Proposition~\ref{prop:counterexample-reduction}
would produce a nonempty finite oriented graph with positive indegrees
satisfying~\eqref{eq:impossible-concavity} for every nonempty proper vertex
set.  Proposition~\ref{prop:no-concavity} rules out such a graph.
\end{proof}
